\documentclass{article}
\usepackage{amsmath}
\def\filler{Text that runs on for a while, so that one line of source fills more than a page. }
\def\fillers{\filler\filler\filler\filler\filler\filler\filler\filler\filler\filler}
\begin{document}
\section{Introduction}\label{intro}
A first paragraph, with some words in it,
and a second line of it, with $x^2 + y^2 = z^2$ inline,
and a third line ending it.

An input file's paragraph,
on two lines.


\begin{itemize}
  \item One item, \emph{emphasised}.
  \item Another item\footnote{With a footnote.}.
\end{itemize}

\begin{equation}
  E = mc^2 \label{eq}
\end{equation}

\begin{center}
\begin{tabular}{l|r}
  Left & Right \\
  \hline
  1 & 2 \\
\end{tabular}
\end{center}

Kerns and glue\hfill at the right\kern 2em end.\hspace{3em}x

\fillers\fillers\fillers\fillers

See Section~\ref{intro} and \eqref{eq}. \mbox{A box} and \fbox{a frame}.
A form: \setbox0\hbox{saved}\pdfxform0 \pdfrefxform\pdflastxform\ and on.
\newpage
Last page, \textbf{bold} and \textit{italic}.
\end{document}
